% Options for packages loaded elsewhere
\PassOptionsToPackage{unicode}{hyperref}
\PassOptionsToPackage{hyphens}{url}
\PassOptionsToPackage{dvipsnames,svgnames,x11names}{xcolor}
%
\documentclass[
  11pt,
  a4paper,
]{ltjsarticle}
\usepackage{amsmath,amssymb}
\usepackage{iftex}
\ifPDFTeX
  \usepackage[T1]{fontenc}
  \usepackage[utf8]{inputenc}
  \usepackage{textcomp} % provide euro and other symbols
\else % if luatex or xetex
  \usepackage{unicode-math} % this also loads fontspec
  \defaultfontfeatures{Scale=MatchLowercase}
  \defaultfontfeatures[\rmfamily]{Ligatures=TeX,Scale=1}
\fi
\usepackage{lmodern}
\ifPDFTeX\else
  % xetex/luatex font selection
\fi
% Use upquote if available, for straight quotes in verbatim environments
\IfFileExists{upquote.sty}{\usepackage{upquote}}{}
\IfFileExists{microtype.sty}{% use microtype if available
  \usepackage[]{microtype}
  \UseMicrotypeSet[protrusion]{basicmath} % disable protrusion for tt fonts
}{}
\makeatletter
\@ifundefined{KOMAClassName}{% if non-KOMA class
  \IfFileExists{parskip.sty}{%
    \usepackage{parskip}
  }{% else
    \setlength{\parindent}{0pt}
    \setlength{\parskip}{6pt plus 2pt minus 1pt}}
}{% if KOMA class
  \KOMAoptions{parskip=half}}
\makeatother
\usepackage{xcolor}
\usepackage[margin=24mm]{geometry}
\setlength{\emergencystretch}{3em} % prevent overfull lines
\providecommand{\tightlist}{%
  \setlength{\itemsep}{0pt}\setlength{\parskip}{0pt}}
\setcounter{secnumdepth}{5}
\ifLuaTeX
  \usepackage{selnolig}  % disable illegal ligatures
\fi
\IfFileExists{bookmark.sty}{\usepackage{bookmark}}{\usepackage{hyperref}}
\IfFileExists{xurl.sty}{\usepackage{xurl}}{} % add URL line breaks if available
\urlstyle{same}
\hypersetup{
  colorlinks=true,
  linkcolor={blue},
  filecolor={Maroon},
  citecolor={Blue},
  urlcolor={blue},
  pdfcreator={LaTeX via pandoc}}

\author{}
\date{}

\begin{document}

\hypertarget{ux7279ux7570ux5024ux5206ux89e3svd}{%
\section{特異値分解（SVD）}\label{ux7279ux7570ux5024ux5206ux89e3svd}}

任意の実行列

\[
A\in\mathbb{R}^{m\times n}
\]

は

\[
\boxed{A=U\Sigma V^\top}
\]

と分解できます。

\hypertarget{ux4f55ux3092ux3057ux3066ux3044ux308bux306eux304b}{%
\subsection{1.
何をしているのか}\label{ux4f55ux3092ux3057ux3066ux3044ux308bux306eux304b}}

ベクトル \texttt{x} に \texttt{A} を作用させる過程を

\[
x\xrightarrow{V^\top}V^\top x\xrightarrow{\Sigma}\Sigma V^\top x\xrightarrow{U}Ax
\]

と見ます。

\begin{itemize}
\tightlist
\item
  \texttt{V\^{}T}: 入力空間を右特異ベクトル基底へ回転・鏡映する
\item
  \texttt{Σ}: 各直交方向を \texttt{σ\_i} 倍する
\item
  \texttt{U}: 出力空間の左特異ベクトル方向へ回転・鏡映する
\end{itemize}

つまり一般の線形写像を「直交変換 → 軸ごとの伸縮 →
直交変換」に分けています。

\hypertarget{ux7279ux7570ux5024ux306fux3069ux3046ux6c42ux307eux308bux304b}{%
\subsection{2.
特異値はどう求まるか}\label{ux7279ux7570ux5024ux306fux3069ux3046ux6c42ux307eux308bux304b}}

\[
A^\top A=V\Sigma^\top\Sigma V^\top
\]

なので、\texttt{V} の列は \texttt{A\^{}TA} の固有ベクトルです。また

\[
\lambda_i(A^\top A)=\sigma_i^2
\]

です。

同様に \texttt{U} の列は \texttt{AA\^{}T} の固有ベクトルです。

したがって

\[
Av_i=\sigma_i u_i
\]

となります。右特異ベクトル方向 \texttt{v\_i} の単位ベクトルを \texttt{A}
が \texttt{σ\_i} 倍し、左特異ベクトル方向 \texttt{u\_i}
へ送るという意味です。

\hypertarget{ux30e9ux30f3ux30afux3068ux306eux95a2ux4fc2}{%
\subsection{3.
ランクとの関係}\label{ux30e9ux30f3ux30afux3068ux306eux95a2ux4fc2}}

非零特異値の個数がランクです。

\[
\operatorname{rank}(A)=\#\{i:\sigma_i>0\}
\]

零特異値に対応する右特異ベクトルは \texttt{ker(A)} を張ります。

\hypertarget{ux5916ux7a4dux306eux548cux3068ux3057ux3066ux898bux308b}{%
\subsection{4.
外積の和として見る}\label{ux5916ux7a4dux306eux548cux3068ux3057ux3066ux898bux308b}}

SVD は

\[
A=\sum_{i=1}^r\sigma_i u_i v_i^\top
\]

とも書けます。

各 \texttt{u\_i\ v\_i\^{}T} はランク1行列です。したがって \texttt{A}
は、重要度 \texttt{σ\_i} を持つランク1成分の和として理解できます。

この見方から画像圧縮や低ランク近似が自然に出てきます。

\hypertarget{ux6761ux4ef6ux6570ux3068ux306eux95a2ux4fc2}{%
\subsection{5.
条件数との関係}\label{ux6761ux4ef6ux6570ux3068ux306eux95a2ux4fc2}}

full rank の場合

\[
\kappa_2(A)=\frac{\sigma_{\max}}{\sigma_{\min}}
\]

です。最小特異値が小さいとは、ある入力方向がほとんど潰されることを意味します。その方向を逆に復元しようとすると誤差が大幅に増幅されます。

\hypertarget{svd-ux304cux7d71ux4e00ux3059ux308bux3082ux306e}{%
\subsection{6. SVD
が統一するもの}\label{svd-ux304cux7d71ux4e00ux3059ux308bux3082ux306e}}

SVD から直接

\begin{itemize}
\tightlist
\item
  ランク
\item
  四つの基本部分空間
\item
  擬似逆行列
\item
  最小二乗
\item
  条件数
\item
  低ランク近似
\item
  PCA
\item
  画像圧縮
\item
  ノイズ除去
\end{itemize}

を説明できます。

SVD
を単なる計算手法ではなく、\textbf{線形写像の幾何学的な標準形}として理解することが最重要です。

\end{document}
